\subsection{Finite-bit realization}\label{sec:bits}

The recursive description contains known real probabilities and one
countable distribution. This section gives finite algorithms for both.
All interval endpoints used by the implementation are rational.

\begin{lemma}[Exact comparison with a computable probability]\label{lem:lazy}
Suppose an oracle gives a rational interval containing
\(p\in[0,1]\) of width at most \(2^{-m}\) at stage \(m\), in
\(O((B+m)^c)\) bit operations. A fair-bit algorithm samples
\(\Bern(p)\) exactly with \(O(B^c+1)\) expected bit work.
It terminates almost surely, including when \(p=0\), \(p=1\), or
\(p\) is dyadic.
\end{lemma}
\paragraph{Proof idea.}
Compare one uniform binary expansion with successively narrower certified
intervals for the probability. A comparison remains undecided only when the
uniform value lies close to the target, so its geometric probability tail
pays for precision refinement.

\begin{proof}
Reveal a uniform binary expansion \(U\), maintaining its length-\(m\)
dyadic interval. Stop when the two intervals are disjoint, and output
whether \(U<p\). If no decision has been made at stage \(m\), then
\(|U-p|\le2^{1-m}\). This event has probability at most
\(4\,2^{-m}\). The equality event has probability zero.
Summing the stage costs against this geometric tail proves the
expectation bound. Rational choices among finitely many categories
are handled either by an integer rejection draw or by successive
conditional coins.
\end{proof}

On bounded real arguments, certified Taylor sums with rounding
after each operation give tanh intervals of width \(2^{-m}\)
in \(O((B+m)^4)\) bit operations. Section~\ref{sec:fixedprecision}
gives the error analysis, first for the less immediate countable
distribution and then for the elementary functions.

For a nonzero row, \(\rho=\sigma r_{\ell-1}\) is a positive
rational of \(O(\bits)\) bits. In particular,
\(\rho\ge2^{-O(\bits)}\). Dividing by \(\tanh\rho\), \(r_\ell\),
or \(R\) requires only \(O(\bits)\) guard bits: on the bounded
argument range, \(\tanh\rho\) and \(R\) are at least a constant
multiple of \(\rho\). For \(R\), use
\(R=\frac12\int_{a-\rho}^{a+\rho}\sech^2u\,du\) and
\(|a|+\rho\le5\).
Zero denominators are avoided by the explicit zero-row branch.

\subsubsection{An exact arity draw}

For \(0<\rho\le2\), positivity gives
\[
F_\rho(5/4)=\sec^2(\rho/2)\le4,\qquad
d_k\le4(4/5)^k.
\]
Here \(\cos1\ge1/2\).
Also \(\rho/\tanh\rho\le7/3<3\) and \(C_k\ge1\), so
\begin{equation}\label{eq:arityenvelope}
\pi_k\le12(4/5)^k.
\end{equation}
Propose \(k\ge0\) with probability
\(\gamma_k=(1/5)(4/5)^k\), and accept with probability
\(\pi_k/(60\gamma_k)\). Each trial succeeds with probability
\(1/60\). Its accepted index has law \(\pi\).
No source signs are requested before acceptance.

\subsubsection{A fourth-power bound for scalar bit work}
\label{sec:fixedprecision}

The scalar probabilities can be refined without carrying a growing
common rational denominator.  We give the details because this reduces
the bit overhead in the network theorem.

\begin{lemma}[Fixed-precision coefficient evaluation]
\label{lem:coefficient-four}
Let \(0<\rho\le2\) be a dyadic rational whose description has at
most \(B\) bits.  For integers \(K\ge0\) and \(m\ge1\), the
coefficient \(d_K\) of
\[
  \sech^2(\rho\sqrt{1-z})=\sum_{j\ge0}d_jz^j
\]
has a certified dyadic enclosure of width \(2^{-m}\), computable
in \(O((B+K+m+1)^4)\) bit operations.  The constant is absolute,
and schoolbook multiplication and division suffice.
\end{lemma}

\paragraph{Proof idea.}
Compute the coefficients of the reciprocal denominator first, then invert
its triangular coefficient system. Explicit error bounds allow rounding
after every operation, so the rational denominators never grow with the
number of operations.

\begin{proof}
It suffices to construct an approximation with error \(2^{-m-2}\).
In the following calculation replace the requested precision \(m\)
by \(m+2\) for this purpose.  Write
\[
 A(z)=\cosh^2(\rho\sqrt{1-z})=\sum_{j\ge0}a_jz^j.
\]
The coefficients have alternating signs before the constant term,
so the power series for cosh gives
\begin{equation}\label{eq:coefficient-norm}
 a_0\ge1,\qquad
 \sum_{j\ge0}|a_j|=\cosh^2(\rho\sqrt2)<1024.
\end{equation}
Indeed, \(\rho\sqrt2<3\), \(e<3\), and
\(\cosh^2 3<e^6<729\).  To bound the reciprocal, write
\(\sqrt{1-z}=u+iv\).  On \(|z|\le1\),
\[
 |z|^2=(u^2+v^2-1)^2+4v^2,
 \qquad |v|\le1/2.
\]
Consequently
\[
 |A(z)|=\sinh^2(\rho u)+\cos^2(\rho v)
 \ge\cos^2 1\ge1/4,
\]
and hence
\begin{equation}\label{eq:reciprocal-coefficient-bound}
 |d_j|\le4\qquad(j\ge0)
\end{equation}
by Cauchy's coefficient estimate.

Define integer precision parameters
\begin{align*}
 T&=m+12(K+1)+\lceil\log_2(K+1)\rceil+20,\\
 M&=T+32,\\
 P&=T+8M+10\lceil\log_2(M+1)\rceil+64.
\end{align*}
All operations below round to \(P\) binary fractional places, with
absolute error at most \(2^{-P}\) per rounding.  Their integer
parts have \(O(M)\) bits.  The input \(\rho\) is retained exactly.

For \(0\le j\le M\), compute
\[
 c_j=\frac{(2\rho)^{2j}}{(2j)!},\qquad
 c_0=1,\qquad
 c_{j+1}=c_j\frac{4\rho^2}{(2j+1)(2j+2)}.
\]
Use these values to form the coefficients through degree \(K\) of
\[
 A_M(z)=\frac12\left(1+\sum_{j=0}^M c_j(1-z)^j\right).
\]
The binomial coefficients are exact integers, obtained by Pascal's
recurrence; each has at most \(M+1\) bits.  Terms of degree above
\(K\) may be discarded, but this is unnecessary for the stated
complexity.

The multipliers in the recurrence for \(c_j\) are at most 16.
Its accumulated rounding error is at most
\(16^{M+1}2^{-P}\) at every index.  Expanding the binomials adds
at most a factor \((M+1)2^M\).  The remaining additions and
roundings add at most \((M+1)^2 2^{-P}\) to a coefficient error.
The specified value of \(P\) bounds their sum by
\(2^{-T-4}\).

The truncation error on \(|z|\le1\) is at most
\[
 E_M=\frac12\sum_{j>M}\frac{32^j}{(2j)!}
 \le2^{32-4M}\le2^{-T-4}.
\]
For the middle inequality, \(M\ge16\); successive terms have
ratio below \(1/16\), and
\(34!\ge2^{118}\) follows from
\(16!\ge2^{44}\), \(\prod_{j=17}^{32}j\ge2^{64}\), and
\(33\cdot34\ge2^{10}\).
Cauchy's estimate applies to each coefficient error.
Thus the computed coefficients satisfy
\begin{equation}\label{eq:coefficient-input-error}
 |\widehat a_j-a_j|\le\Delta:=2^{-T},\qquad 0\le j\le K,
 \qquad \widehat a_0\ge1/2.
\end{equation}

Now use the reciprocal recurrence with the same precision:
\[
 \widehat d_0=\operatorname{round}_P(1/\widehat a_0),\qquad
 \widehat d_k=\operatorname{round}_P\left(
 -\frac{\sum_{j=1}^k\widehat a_j\widehat d_{k-j}}
        {\widehat a_0}\right).
\]
Each product is rounded before summation.  Put
\(E_k=\max_{0\le j\le k}|\widehat d_j-d_j|\).
Provided \(E_{k-1}\le1\), equations
\eqref{eq:coefficient-norm}--\eqref{eq:coefficient-input-error}
give
\[
 E_0\le3\Delta,\qquad
 E_k\le2048E_{k-1}+(12K+8193)\Delta\quad(1\le k\le K).
\]
To see the second estimate, the error in the numerator is at most
\(1024E_{k-1}+5k\Delta+k2^{-P}\).  Division by
\(\widehat a_0\) multiplies this by at most two.  Changing
\(1/a_0\) to \(1/\widehat a_0\) contributes at most
\(8192\Delta\), since the exact numerator is bounded by 4096.
The final rounding contributes at most \(2^{-P}\).

Iterating the affine recurrence gives
\[
 E_k\le\left(3+\frac{12K+8193}{2047}\right)2^{11k}\Delta
 \le2^{11k+14}(K+1)\Delta\quad(0\le k\le K).
\]
The choice of \(T\) makes this less than one at every index and
gives \(E_K<2^{-m}\), validating the assumed bound throughout.
The known error bounds supply a certified interval about
\(\widehat d_K\).

There are \(O(M^2+K^2)\) arithmetic operations.  Every operand
has \(O(B+P+M)=O(B+K+m+1)\) bits, including exact binomial
coefficients, products with the input dyadic rational, and rounded
quotients.  Schoolbook arithmetic therefore costs
\(O((B+K+m+1)^4)\) bit operations in total.
\end{proof}

The elementary exponential series on bounded real arguments can
be evaluated by the same fixed-precision method in cubic bit work:
use linearly many terms, linearly many guard bits, and round after
each multiplication or division.  Their total amplification is
bounded by a constant on the argument range used here.  A fourth-power
bound therefore covers tanh, the bias probabilities, and their
rational combinations.  Nonzero denominators require \(O(\bits)\)
guard bits as established above.  The exact integers defining
\(C_K\), and the factor \((5/4)^K\) in the arity acceptance
probability, require \(O(K)\) additional bits.

It follows from Lemma~\ref{lem:coefficient-four} that the accepted
arity can be drawn with \(O(\bits^4)\) expected scalar bit work.
Both the proposal arity and the precision stage have geometric
tails, uniformly in the row and the context.  The mean number of
proposals is 60.  All other scalar work at a reached call obeys the
same bound.  Lemma~\ref{lem:stopping} then gives
\begin{equation}\label{eq:fourth-power-overhead}
 \text{expected online bit work}
 =O\!\left(\bits^4\E N_{\rm calls}\right).
\end{equation}
The preprocessing and storage bounds are unchanged.


\paragraph{Proof idea.}
Preprocessing supplies short rational row descriptions and sampling tables.
At every reached call, the scalar routines have a uniform expected cost.
Predictable charging then sums those costs over the random recursion tree.

\begin{proof}[Completion of Theorem~\ref{thm:upper}]
Every stored parameter has \(O(\bits)\) bits because
\(P=O(\bits)\). A dense center sum costs \(O(n\bits^2)\);
the bounded scalar endpoints cost \(O(\bits^4)\).
The prefix tables and all row data are obtained within
\(O(Dn^2\bits^6)\) time and \(O(Dn^2\bits)\) storage.

A child index is drawn from the absolute integer weights by an
integer rejection draw of success probability at least \(1/2\)
and binary search in the stored prefix sums. This costs
\(O(\bits^2)\) expected bit operations, including the address.
All other scalar work per reached call is
\(O(\bits^4)\) in expectation, uniformly over its context
and the prior history. Lemma~\ref{lem:stopping}, summed over the
recursion, bounds total bit work by
\(O(\bits^4\E N_{\rm calls})\).
No independence between the duration and output of a child is used.
The exactness identities in Sections~\ref{sec:algorithm}--\ref{sec:bias}
and almost sure termination prove the required output law.
\end{proof}

For completeness, a dense exact sampler is available for every fixed
\(s\). Evaluate the network with certified intervals and compare its
output probability to one uniform binary expansion.
To obtain output width \(2^{-m}\), it suffices to use internal
precision
\[
m+O\bigl(b+\log(nD)+D\log^+s\bigr).
\]
This follows by propagating interval errors through the
\(s\)-Lipschitz rows. At \(D=L\), \(b=20L\), Lemma~\ref{lem:lazy}
gives \(\widetilde O_s(n^2)\) expected bit work.
\paragraph{One uniform choice of sampler.}
At \(D=L\), the refined bounds and the dense upper bound can be
attained by one weight-only selection rule. If \(S>2\), choose
the dense algorithm. Otherwise, after finding \(S\),
consider the rational comparison gains \(q=1+k/D\), \(1\le k\le D\),
and the additional gain \(1+2^{-16}\), retaining those at least
\(\max\{1,S\}\). Evaluate the bounds of
Propositions~\ref{prop:fixedrate}
and~\ref{prop:fusedrate} where applicable, together with the uniform recursive bound, the
critical bound when \(S\le1\), and the dense envelope \(n^2\).
Approximate the logarithm of each envelope by a certified upper
value within one, and select the recursive or dense algorithm
corresponding to the smallest value, using the original or fused
factory as required. The common finite-bit factors
are powers of \(\log n\), so this choice achieves the minimum of
the displayed envelopes up to such a factor. The recursive preprocessing is performed once, with reference gain
\(\max\{1,S\}\); the other gains are used only to compare bounds.
The fused factory uses the same stored data and a fixed cutoff.

This selection costs \(O(D\bits^6)\) preprocessing operations.
For completeness, write \(w=u_*^2\) and invert
\[
\Phi(w)=\sqrt w\coth\sqrt w
=\frac{\sum_{j\ge0}w^j/(2j)!}
       {\sum_{j\ge0}w^j/(2j+1)!},\qquad 0\le w\le4.
\]
Its derivative is at least \(1/12\): with \(u=\sqrt w\),
\[
\Phi'(w)=\frac{\sinh(2u)-2u}{4u\sinh^2u}\ge\frac1{12},
\]
using \(\sinh(2u)-2u\ge4u^3/3\) and
\(\sinh u/u\le\sinh2/2<2\); the value at zero is removable.
Certified interval bisection therefore encloses \(w\) to the
required accuracy without an exact equality test. When the midpoint
value interval overlaps the target, the derivative bound encloses
the root in a shorter interval about that midpoint. The smallest
positive candidate gap is at least \(\min\{D^{-1},2^{-16}\}\),
and the fixed-point estimates give
\(w\ge3(q-1)\) and \(1-\lambda\ge(q-1)/2\).
Thus all inverse and logarithmic operations need only \(O(\bits)\)
guard bits. The geometric denominators are bounded below by constant
multiples of the same gap. Logarithms of the large prefactors are
evaluated directly, using stable log-sum formulas and bounds on
small negative exponentials; the prefactors themselves are never
formed.

For any fixed \(1<s\le2\), the candidate
\(q_D=\lceil sD\rceil/D\) is eligible and satisfies
\(s\le q_D<s+1/D\). The fixed-point rates are smooth for
\(s>1\), and their prefactors are locally bounded. Thus replacing
\(s\) by \(q_D\) changes the rate multiplied by \(D\) by
\(O_s(1)\). The same argument applies to the fused rate for
\(s<1+2^{-16}\); the added endpoint covers equality. Consequently
the refined minima with the dense bound hold for this single
uniform algorithm. The selection is completed before the context
is revealed.

The chosen sampler therefore retains its exact output law.
This proves the upper side of Corollary~\ref{cor:threshold}, including
the refined fixed-norm minima.
